\section{More on Access Rights}

\begin{frame}[fragile]{What Are Access Rights?}
	\begin{itemize}
		\item Access rights define model-level permissions in \texttt{ir.model.access}.
		\item For each group + model, you allow or deny:
		\begin{itemize}
			\item \texttt{perm\_read}
			\item \texttt{perm\_write}
			\item \texttt{perm\_create}
			\item \texttt{perm\_unlink}
		\end{itemize}
\begin{block}{Usual File}
\begin{lstlisting}
security/ir.model.access.csv
\end{lstlisting}
\end{block}
	\end{itemize}
\end{frame}

\begin{frame}[fragile]{Access Rights CSV in Detail}
\begin{lstlisting}
id,name,model_id:id,group_id:id,perm_read,perm_write,perm_create,perm_unlink
access_student_user,student.user,model_student_student,school_manager.group_school_user,1,0,0,0
access_student_officer,student.officer,model_student_student,school_manager.group_school_officer,1,1,1,0
access_student_manager,student.manager,model_student_student,school_manager.group_school_manager,1,1,1,1
\end{lstlisting}
\end{frame}

\begin{frame}[fragile]{Breaking Down CSV Columns}
	\begin{itemize}
		\item \textbf{\texttt{id}}: external ID of the access line
		\item \textbf{\texttt{name}}: label
		\item \textbf{\texttt{model\_id:id}}: model XML ID (\texttt{model\_student\_student})
		\item \textbf{\texttt{group\_id:id}}: group XML ID
		\item \textbf{\texttt{perm\_*}}: \texttt{1} allow, \texttt{0} deny
	\end{itemize}
\begin{lstlisting}
model_student_student
\end{lstlisting}
	\begin{itemize}
		\item Odoo creates that XML ID from model \texttt{student.student}.
	\end{itemize}
\end{frame}

\begin{frame}[fragile]{Global Access vs Group Access}
\begin{lstlisting}
id,name,model_id:id,group_id:id,perm_read,perm_write,perm_create,perm_unlink
access_student_all,student.all,model_student_student,,1,0,0,0
\end{lstlisting}
	\begin{itemize}
		\item Empty \texttt{group\_id} means the right applies to everyone.
		\item Use sparingly; prefer explicit group rights.
	\end{itemize}
\end{frame}

\begin{frame}[fragile]{Access Rights for Transient / Wizard Models}
\begin{lstlisting}
access_assign_wizard_user,assign.wizard.user,model_student_assign_wizard,school_manager.group_school_user,1,1,1,1
\end{lstlisting}
	\begin{itemize}
		\item Wizards also need ACL lines.
		\item Users often need create/write/read on transient models to open and submit wizards.
		\item Forgetting wizard ACL is a common ``Access Error'' cause.
	\end{itemize}
\end{frame}

\begin{frame}[fragile]{Testing Access in Practice}
\begin{lstlisting}
# Switch to a non-admin user in shell / tests
Student = self.env['student.student'].with_user(teacher)
Student.check_access_rights('read')
Student.check_access_rights('unlink')  # may raise
\end{lstlisting}
	\begin{itemize}
		\item Always verify rights with a normal user, not only admin.
		\item Admin/superuser can hide ACL mistakes.
	\end{itemize}
\end{frame}

\begin{frame}{Access Rights Best Practices}
	\begin{itemize}
		\item Give least privilege needed for each role.
		\item Managers get unlink; normal users usually do not.
		\item Load security XML groups before the ACL CSV.
		\item Add ACL for every custom model, including wizards.
		\item Document your role matrix (User/Officer/Manager $\times$ CRUD).
	\end{itemize}
\end{frame}
